\chapter{Limitations and Future Directions}\label{chapter-6-limitations-and-publication-roadmap}

\section{Methodological choices}\label{methodological-rigor-and-validations}

A few choices were made to keep the evaluation sound:

\begin{itemize}
\tightlist
\item
  \textbf{Statistical power.} Instead of a single sign test, the main comparison is run over a medium-scale set (8 problems $\times$ 6 seeds) and evaluated with a Wilcoxon signed-rank test, so the primary effect is reported with a significance test rather than as a directional signal (Section~4.3).
\item
  \textbf{Compute matching.} A matched-budget comparison and a compute-to-correct frontier (Section~4.4) let the gain be attributed to which trajectory is distilled rather than to drawing more samples.
\item
  \textbf{Reasoning traces on code.} Running the code arm with thinking on provides the reasoning traces needed to measure epistemic verbalization on code (Section~4.7), not only in the math pilot.
\item
  \textbf{Exact filtering.} The verifier--judge filter is exact (Section~4.6): the judge is never bypassed, so only correct, independent trajectories are distilled.
\item
  \textbf{Boundary characterization.} Running the same pipeline on code and math separates dense execution feedback from value-only feedback, which is what exposes the boundary of the method.
\end{itemize}

\section{Scope and limitations}\label{scope-and-remaining-limitations}

Several limitations remain:

\begin{itemize}
\tightlist
\item
  \textbf{Model and scale.} All results come from a single 4B Qwen model at a personal compute scale, so they are an empirical characterization within one model family rather than a scaling law or an industry-scale benchmark.
\item
  \textbf{Sparse math feedback.} The math pilot is limited by the AIME benchmark, which provides only integer answers with no worked solutions. Without a step-by-step reference, the teacher has no procedure to distill, which is why the model does not escape the flat-reward trap on beyond-capability problems.
\item
  \textbf{Margin sensitivity to base capability.} On stronger base models (e.g.\ 8B and beyond), the unconditioned student explores better on its own, so the teacher-first margin is expected to narrow on reachable tasks; the method is most useful where the base model would otherwise stall.
\item
  \textbf{Generality of the epistemic result.} The suppression measurements use one model family and a fixed marker set, so their generality across tokenizers, scales, and reasoning styles is untested. The beneficial reading of suppression on code (consolidation rather than copying) is a co-occurrence interpreted mechanistically; establishing the reference type as the cause would need a same-domain reference-type ablation, which this study does not run.
\end{itemize}
